\section{Données et courbes de marché}
\label{sec:donnees}

\subsection{Données}

Les données de marché sont celles de Bloomberg à la clôture du 8 septembre 2026, pour un règlement le
10 septembre. Cette date de règlement est l'origine des temps, $t=0$ ; le temps est mesuré en années sur une base
ACT/365, sans calendrier de jours ouvrés. Le tableau~\ref{tab:donnees} décrit les quatre sources utilisées.

\begin{table}[htbp]
\centering\small
\caption{Données de marché (Bloomberg, clôture du 8 septembre 2026)}
\label{tab:donnees}
\begin{tabularx}{\textwidth}{@{}l X X@{}}
\toprule
Source & Contenu & Usage \\
\midrule
Courbe €STR & 32 taux de swaps OIS €STR, de 1 semaine à 50 ans & courbe sans risque \\
Cube de swaptions & volatilités normales des swaptions EUR à la monnaie, 21 expiries $\times$ 14 tenors & calibration du facteur taux \\
OAT & 19 titres d'État (5 zéro-coupons courts, 14 OAT de 2028 à 2072) : coupon, maturité, prix clean & courbe OAT, sous-jacent \\
Historiques & rendements OAT génériques et taux de swaps €STR à 2, 5, 10, 15, 20, 25 et 30 ans, quotidiens & estimation du facteur spread \\
\bottomrule
\end{tabularx}
\end{table}

Les historiques couvrent la période du 1\textsuperscript{er} janvier 2021 au 8 septembre 2026, soit 1\,475 jours
ouvrés. Ils sont exportés avec des dates mal formées, dates Excel et texte au format américain mêlés ; la grille
étant quotidienne et complète, on la reconstruit. Les séries s'arrêtent à la date de valorisation, de sorte
qu'aucune information postérieure n'entre dans l'estimation.

S'y ajoute une courbe de spread repo, écart entre le taux de pension livrée de l'OAT et le taux zéro €STR de même
durée. Faute de cotations accessibles, \emph{cette courbe est fictive} : elle vaut +20~pb à trois semaines, puis
croît linéairement jusqu'à +75~pb à dix ans. Les options étudiées ont une échéance de dix ans au plus, de sorte
que la courbe couvre toutes les échéances sans jamais être prolongée. Une courbe de marché la remplacerait au même
format, sans modification du code ; la section~\ref{sec:resultats} mesure l'effet de son niveau.

\begin{figure}[htbp]
\centering
\includegraphics[width=\textwidth]{donnees}
\caption{Données au 8 septembre 2026. À gauche, taux de swaps €STR et rendements des OAT ; au centre,
volatilités normales des swaptions à la monnaie ; à droite, historique du spread OAT–€STR à quatre tenors.}
\label{fig:donnees}
\end{figure}

La figure~\ref{fig:donnees} résume ces données. Aux tenors longs, le spread OAT–€STR s'est élargi d'une
centaine de points de base depuis 2021, avec deux sauts nets, en juin 2024 puis à l'automne 2024. À deux ans, il est resté négatif presque
tout au long de 2021 et 2022.

\subsection{Courbe sans risque €STR}

Les swaps OIS €STR, qui ont remplacé les swaps EONIA \citep[CM1, p.~12]{Hillairet2026}, servent à construire la
courbe sans risque $P^M(0,T)$. Le taux OIS est le rendement du seul taux court au jour le jour
\citep[cours~5, éq.~II.17]{MPR2025}, et la même courbe actualise les flux collatéralisés
\citep[section~3.4]{Chaix2021}. Un swap de maturité $T$ et de taux fixe $S$, à jambe fixe annuelle, vaut zéro à
l'initiation :
\begin{equation}
S \sum_{i} \tau_i\,P^M(0,T_i) + P^M(0,T) = 1, \qquad T_i = T, T-1, \dots
\end{equation}
On résout cette équation pilier après pilier (\emph{bootstrap}), $P^M(0,T)$ étant la seule inconnue une fois
les facteurs d'actualisation intermédiaires interpolés linéairement en taux zéro. La courbe obtenue reprice
les 32 swaps à 0,25~pb près.

Le modèle utilise le forward instantané $f^M(0,t) = -\partial_t \log P^M(0,t)$, qui entre dans le terme de
recalage. Une interpolation linéaire des taux zéro donnerait un forward discontinu ; on interpole donc
$\log P^M(0,T)$ par une spline cubique, ce qui rend $f^M$ continu et dérivable. La courbe €STR étant
décroissante au-delà de vingt ans, le forward long y est nettement inférieur aux taux zéro : environ 1,4~\%
à cinquante ans, pour un taux zéro proche de 2,8~\%.

\subsection{Courbe OAT et sous-jacent}

La courbe $\bar P^M(0,T)$ est la courbe d'actualisation qui reprice les OAT. Un bootstrap titre par titre
propagerait à toute la courbe l'erreur de prix du titre le plus court et dépendrait des particularités de chaque
titre \citep[CM1, p.~27]{Hillairet2026}, alors que le modèle a besoin d'un forward régulier. On ajuste donc la
paramétrisation de \citet{NelsonSiegel1987} étendue par \citet{Svensson1994}, qu'utilisent la plupart des
banques centrales \citep[cours~3, slide~109]{MPR2025}, sur le taux zéro continu :
\begin{equation}
\begin{split}
\bar R(T) = \beta_0 + \beta_1\,\frac{1-e^{-T/\lambda_1}}{T/\lambda_1}
&+ \beta_2\Big(\frac{1-e^{-T/\lambda_1}}{T/\lambda_1} - e^{-T/\lambda_1}\Big)
+ \beta_3\Big(\frac{1-e^{-T/\lambda_2}}{T/\lambda_2} - e^{-T/\lambda_2}\Big), \\
\bar P^M(0,T) &= e^{-\bar R(T)\,T}.
\end{split}
\end{equation}
Les six paramètres sont ajustés par moindres carrés sur les prix dirty (coupon couru inclus) des dix-neuf
titres. Chaque erreur de prix est divisée par la duration du titre, ce qui la rend homogène à une erreur de
rendement ; sans cette pondération, les titres longs domineraient l'ajustement. Le coupon couru, calculé en
ACT/ACT, sert à passer du prix clean (pied de coupon) au prix dirty, aujourd'hui comme à l'échéance d'une
option.

L'erreur quadratique moyenne est de 4,1~pb de rendement (annexe~C). Les résidus se concentrent sur l'extrême long
terme, que six paramètres ne peuvent pas reproduire : les OAT 2050 (coupon de 1,5~\%) et 2072 (0,5~\%) se
traitent cher, par demande de duration et de convexité, et l'OAT 2057 sous la courbe (résidu de $-14{,}5$~pb). Ces
écarts pèsent peu sur le prix de l'option. Son forward est calculé à partir du prix de marché du sous-jacent
(section~\ref{sec:evaluation}), et la courbe OAT n'y entre qu'à travers la répartition des flux et la loi du
spread.

\begin{figure}[htbp]
\centering
\includegraphics[width=\textwidth]{courbe_oat}
\caption{Courbes zéro OAT (Nelson-Siegel-Svensson) et €STR, et spread OAT–€STR en taux zéro et en taux forward
instantané.}
\label{fig:courbe_oat}
\end{figure}

Le spread zéro $\bar R(T) - R(T)$ et le spread forward $\bar f^M - f^M$ de la figure~\ref{fig:courbe_oat}
donnent la partie déterministe du facteur spread. Le spread court est faible : le taux OAT contient, outre le
risque de défaut, des primes de liquidité et de collatéral \citep[p.~8]{Russo}, et le spread peut changer de
signe, comme à deux ans en 2021-2022. Un modèle gaussien, où le spread n'est pas contraint à rester positif, est
donc admissible.

Les options portent sur l'OAT 4,10~\% du 25 mai 2046 (FR0014015MU5), de maturité résiduelle 19,7 ans, dont le
prix clean est 91,690 et le prix dirty 92,903. La courbe la reprice à 2,0~pb de rendement près. Les échéances
des options sont 1, 2, 5, 7 et 10 ans ; à dix ans, il reste encore dix flux au titre.
